\documentclass[12pt,a4paper]{article}

\usepackage[spanish]{babel}
\usepackage[utf8]{inputenc}
\usepackage[T1]{fontenc}
\usepackage{lmodern}
\usepackage{geometry}
\usepackage{setspace}
\usepackage{enumitem}
\usepackage{booktabs}
\usepackage{xcolor}
\usepackage{hyperref}

\geometry{margin=2.5cm}
\setstretch{1.15}
\hypersetup{colorlinks=true,linkcolor=black,urlcolor=blue,citecolor=black}

\title{\textbf{Reporte Técnico UX/UI para Adaptación Responsive en Teléfonos}\\\large CV Builder (GitHub Pages)}
\author{DanielRestrepoGaleano/cv-template}
\date{06 de mayo de 2026}

\begin{document}
\maketitle
\thispagestyle{empty}

\vspace{0.5cm}
\noindent\textbf{Objetivo del documento:} presentar, sin ejecutar cambios de código, el conjunto de ajustes necesarios de diseño UX/UI para que la aplicación web se adapte correctamente a teléfonos móviles, manteniendo su naturaleza estática sobre GitHub Pages.

\vspace{0.8cm}
\tableofcontents
\newpage

\section{Resumen ejecutivo}
La aplicación cuenta con una base funcional sólida para escritorio y capacidades relevantes (editor, vista previa, exportaciones y módulos de IA), pero actualmente presenta una adaptación móvil parcial. El panel de chat ya tiene reglas responsive específicas, mientras que el layout principal (toolbar, editor lateral, vista de CV y densidad tipográfica) no está optimizado integralmente para pantallas pequeñas.

Para lograr una experiencia móvil consistente, se requiere una estrategia \textit{mobile-first} que reorganice navegación, jerarquía visual, entradas de formulario y estructura de lectura del CV. Este reporte prioriza cambios de UX/UI en fases, define criterios de aceptación y propone una matriz de validación orientada a GitHub Pages (sitio estático, sin backend).

\section{Contexto y alcance}
\subsection{Contexto técnico}
\begin{itemize}[leftmargin=1.2cm]
  \item Sitio estático desplegable en GitHub Pages.
  \item Frontend en HTML, CSS y JavaScript vanilla.
  \item Sin servidor propio para renderizado dinámico de interfaces.
\end{itemize}

\subsection{Alcance del reporte}
\begin{itemize}[leftmargin=1.2cm]
  \item Identificar brechas actuales de adaptabilidad móvil.
  \item Definir cambios necesarios de UX/UI (sin implementación).
  \item Entregar backlog priorizado para ejecución futura.
\end{itemize}

\subsection{Fuera de alcance}
\begin{itemize}[leftmargin=1.2cm]
  \item Cambios de código (HTML/CSS/JS).
  \item Refactor de arquitectura de datos.
  \item Cambios de hosting o infraestructura.
\end{itemize}

\section{Diagnóstico del estado actual (enfoque móvil)}
\subsection{Hallazgos principales}
\begin{enumerate}[leftmargin=1.2cm]
  \item \textbf{Layout principal no mobile-first:} la app mantiene un esquema de escritorio con panel lateral fijo y vista principal de CV, sin una transición estructural completa para teléfonos.
  \item \textbf{Toolbar densa para pantallas pequeñas:} alta cantidad de acciones en la cabecera, con riesgo de saturación visual, wraps incómodos y dificultad de toque.
  \item \textbf{Vista previa de CV con estructura compleja para móvil:} composición tipo dos columnas y tipografías pensadas para escritorio; la lectura en ancho reducido pierde escaneo rápido.
  \item \textbf{Formularios extensos con alta carga cognitiva:} en móvil faltan patrones de progresión por bloques (pasos/acordeones/secciones compactas) para edición eficiente.
  \item \textbf{Cobertura responsive parcial:} el módulo de chat sí contempla breakpoints, pero el resto de la interfaz no presenta el mismo nivel de adaptación.
\end{enumerate}

\subsection{Impacto UX estimado}
\begin{itemize}[leftmargin=1.2cm]
  \item Mayor tiempo para completar/editar CV desde teléfono.
  \item Incremento en errores de interacción (toques imprecisos, scroll excesivo).
  \item Reducción de legibilidad y percepción de calidad visual en móvil.
\end{itemize}

\section{Objetivos UX/UI para móvil}
\begin{itemize}[leftmargin=1.2cm]
  \item Garantizar usabilidad completa en anchos de 320--430 px.
  \item Simplificar navegación y priorizar acciones clave.
  \item Mejorar legibilidad, ritmo visual y accesibilidad táctil.
  \item Mantener consistencia con experiencia escritorio y exportaciones.
\end{itemize}

\section{Cambios necesarios de diseño UX/UI}
\subsection{1) Navegación y arquitectura de información}
\begin{itemize}[leftmargin=1.2cm]
  \item Reemplazar la toolbar saturada por una navegación en capas: acciones primarias visibles y acciones secundarias en menú contextual.
  \item Definir una jerarquía de acciones móvil: \textit{Editar}, \textit{Vista previa}, \textit{Exportar} y \textit{IA}.
  \item Incluir un patrón de navegación persistente de bajo costo cognitivo (por ejemplo, barra inferior o tabs compactas).
\end{itemize}

\subsection{2) Estructura del layout principal}
\begin{itemize}[leftmargin=1.2cm]
  \item Migrar de layout lateral de escritorio a flujo vertical móvil.
  \item Evitar presentar editor y preview simultáneamente en pantallas pequeñas; usar conmutación clara entre vistas.
  \item Reducir dependencias de altura fija para prevenir conflictos con barras del navegador móvil.
\end{itemize}

\subsection{3) Experiencia de edición en formularios}
\begin{itemize}[leftmargin=1.2cm]
  \item Introducir secciones colapsables para disminuir longitud percibida.
  \item Estandarizar tamaños táctiles (alto mínimo de controles y botones) para interacción cómoda con pulgar.
  \item Mejorar separación visual entre grupos de campos y acciones de agregar/eliminar ítems repetibles.
\end{itemize}

\subsection{4) Vista previa del CV en teléfono}
\begin{itemize}[leftmargin=1.2cm]
  \item Definir versión de lectura móvil con prioridad de contenido, no réplica literal del lienzo de escritorio.
  \item Revisar estructura de columnas para lectura lineal en móvil.
  \item Ajustar escala tipográfica y espaciados para mantener contraste y escaneo rápido.
\end{itemize}

\subsection{5) Modales y componentes flotantes}
\begin{itemize}[leftmargin=1.2cm]
  \item Unificar comportamiento de modales en móvil (alto máximo, scroll interno, botón de cierre siempre visible).
  \item Validar que paneles flotantes (como chat) no bloqueen acciones críticas ni campos en foco.
  \item Asegurar consistencia de patrones de interacción entre módulos (IA import, chat settings, ATS).
\end{itemize}

\subsection{6) Accesibilidad y ergonomía móvil}
\begin{itemize}[leftmargin=1.2cm]
  \item Asegurar objetivos táctiles mínimos y separación adecuada entre controles.
  \item Garantizar contraste visual y legibilidad de etiquetas, estados y mensajes.
  \item Mejorar estados de foco/teclado para inputs en contextos de pantalla reducida.
\end{itemize}

\subsection{7) Rendimiento percibido en GitHub Pages}
\begin{itemize}[leftmargin=1.2cm]
  \item Priorizar carga inicial ligera para red móvil.
  \item Minimizar bloqueos de interfaz por procesos no esenciales en primer render.
  \item Preservar un flujo usable incluso con latencia alta o dispositivos de gama media.
\end{itemize}

\section{Backlog priorizado de implementación futura}
\subsection{Prioridad alta (P0)}
\begin{enumerate}[leftmargin=1.2cm]
  \item Rediseño de navegación móvil y simplificación de toolbar.
  \item Reestructuración del layout principal para flujo vertical con conmutación editor/preview.
  \item Ajuste de legibilidad base (tipografía, espaciados, densidad) para anchos de teléfono.
\end{enumerate}

\subsection{Prioridad media (P1)}
\begin{enumerate}[leftmargin=1.2cm]
  \item Reorganización de formularios con secciones colapsables.
  \item Optimización de modales para pantallas pequeñas.
  \item Homologación de patrones entre módulos de IA y ATS.
\end{enumerate}

\subsection{Prioridad baja (P2)}
\begin{enumerate}[leftmargin=1.2cm]
  \item Microinteracciones y refinamientos visuales.
  \item Mejora incremental de rendimiento percibido y estados vacíos/mensajes.
\end{enumerate}

\section{Criterios de aceptación UX/UI}
Se considera que la aplicación está adaptada a teléfono cuando:
\begin{itemize}[leftmargin=1.2cm]
  \item Todas las tareas críticas (editar, revisar, exportar) pueden completarse con una mano y sin zoom.
  \item No hay desbordes horizontales en 320 px, 360 px y 390 px.
  \item El tiempo de localización de acciones primarias disminuye de forma perceptible frente al estado actual.
  \item La lectura del CV en móvil mantiene claridad de jerarquía y continuidad narrativa.
\end{itemize}

\section{Matriz mínima de validación sugerida}
\begin{center}
\begin{tabular}{@{}p{3.2cm}p{4.2cm}p{6.6cm}@{}}
\toprule
\textbf{Dispositivo} & \textbf{Navegador} & \textbf{Escenarios a validar} \\
\midrule
320x568 & Chrome Android & Navegación principal, edición de campos largos, apertura/cierre de modales \\
360x800 & Chrome Android & Flujo completo: editar $\rightarrow$ preview $\rightarrow$ exportar \\
390x844 & Safari iOS & Teclado móvil, foco en inputs, paneles flotantes y legibilidad general \\
412x915 & Chrome Android & Formularios con múltiples ítems y comportamiento de scroll \\
\bottomrule
\end{tabular}
\end{center}

\section{Riesgos y mitigaciones}
\begin{itemize}[leftmargin=1.2cm]
  \item \textbf{Riesgo:} introducir inconsistencia entre escritorio y móvil.\\
  \textbf{Mitigación:} definir tokens visuales y reglas de diseño compartidas.
  \item \textbf{Riesgo:} degradar experiencia de exportación al cambiar jerarquías visuales.\\
  \textbf{Mitigación:} separar claramente vista de edición móvil y estilos de salida (PDF/Word/ATS).
  \item \textbf{Riesgo:} complejidad de interacción por exceso de acciones simultáneas.\\
  \textbf{Mitigación:} priorización funcional estricta y progressive disclosure.
\end{itemize}

\section{Conclusión}
La aplicación puede adaptarse exitosamente a teléfonos sin abandonar su modelo estático en GitHub Pages. El principal trabajo requerido no es de infraestructura, sino de rediseño UX/UI orientado a móvil: simplificar navegación, reorganizar layout, mejorar legibilidad y estandarizar componentes interactivos. Con el backlog propuesto, se puede ejecutar una transición ordenada y medible hacia una experiencia móvil profesional.

\vspace{0.7cm}
\noindent\textbf{Nota:} este entregable corresponde exclusivamente a un reporte técnico. No se realizaron cambios funcionales en la aplicación.

\end{document}
